\section{Character smoothing from half-deck information}
\label{frames:char:section}

We now use half-deck information to control the Fourier matrix of the
whole permutation. Fix one coordinate half-partition for a block, let
$P$ denote a long segment of coordinate layers, and let $B$ be an
independent pass within each half. For a unitary representation $\rho$,
write $F=\mathbb E\rho(B)$. There are two useful orders for these segments:
\[
\begin{array}{c|c|c}
\text{block}&\text{Fourier matrix}&
       \text{positive form obtained by convexity}\\
\hline
BP&F\,\mathbb E\rho(P)&\rho(P)^*F^*F\rho(P)\\
PB&(\mathbb E\rho(P))F&\rho(P)FF^*\rho(P)^*
\end{array}
\]
In the first order, the pass acts on fixed output halves of $P$.
The conjugated form can be controlled by preimages of those halves.
In the second order, the pass acts on fixed input halves of $P$.
We will repair their image laws, then average the resulting subgroup
coordinates one at a time. The two coordinates may remain dependent.
We first establish the common half-list and one-pass inputs, and then
carry out these two conversions.

\begin{theorem}\label{frames:char:main}
Let $n=2^d$, and count one physical Thorp shuffle as one time step, using
$n/2$ independent fair switches. For all sufficiently large integers $d$,
the worst-initial-deck total-variation mixing time at threshold $1/4$
satisfies
\[
 d\le t_{\mathrm{mix}}(d)\le 6060000d.
\]
In fact, the worst-initial-deck distance at time
$60000(101d-1)$ tends to zero as $d\to\infty$.
\end{theorem}

\subsection{The shared half-list input}
\label{frames:char:marginals}

We retain the convention that permutations send initial labels to
positions and that later layers multiply on the left. Write $C_v$ for
the matching subgroup in direction $v\ne0$. Removing the deterministic
coordinate rotation from the physical shuffle gives independent uniform
layers in $C_{b_t}$, where $b_1,b_2,\ldots$ cycles through the coordinate
basis. A block may begin at any phase of this cycle.

\begin{lemma}\label{frames:char:half-list}
Let $\pi$ be the product of $L=100d$ consecutive layers in a periodic
coordinate order on $n=2^d$ positions. Uniformly over ordered lists of
$q\le n/2$ distinct inputs, the law of their images under $\pi$ has
total-variation distance $o(1)$ from the uniform ordered injection.
The same assertion holds for $\pi^{-1}$. Both assertions are uniform
in the starting phase and the order of the coordinate axes.
\end{lemma}

\begin{proof}
Proposition~\ref{shear:prop:fixed-marginal} covers every periodic
ordering of the standard axes, including every starting phase, and
the inverse product obtained by reversing independent involutive
layers. With its omitted-block parameter equal to $2$ and time
$T=100d$, it gives, for $\eta=\pi$ or $\eta=\pi^{-1}$,
\[
 \|\mathcal L(\eta(c_1),\ldots,\eta(c_q))-\pi_q\|_{\mathrm{TV}}
 \le\frac{n^{3/2}}2
       \left(e^{-99d/8}+2e^{-n/16}\right)^{1/2}=o(1),
\]
where $\pi_q$ is the uniform ordered injection law. The fixed-list
estimate uses the hidden-card vector from the chronological field
\eqref{eq:marginal-filtration}.
The terminal comparison uses
\eqref{eq:marginal-terminal-identification} with $\Lambda=\mathbf v$,
disclosing the complete schedule and only the complete predecessor
paths.

For the independent shear realization of the coordinate chain,
\eqref{shear:eq:inverse-coupling} gives the local block identity
\begin{equation}\label{frames:char:conjugated-layers}
 B_0=I,\qquad M_t=B_t^{-1}Q_tB_{t-1},\qquad
 \pi=Q_L\cdots Q_1=B_LM_L\cdots M_1.
\end{equation}
Here $Q_t$ are the original coordinate layers and $M_t$ the virtual
layers coupled to them. Conditional on the independent shears, the
$M_t$ are independent uniform matching switches. The known terminal
$B_L$ relabels only final positions, so the specified input list is
unchanged. Equation~\eqref{shear:eq:terminal-factorization} identifies
its virtual terminal posterior with its chronological vector after all
shears are disclosed. The independent shear realization
starts with $B_0=I$;
the general first-basis realization retains the input map
$L_0^{-1}$ from \eqref{eq:frames:tuple-input-map}.
\end{proof}

\subsection{The character input and the trace of one pass}
\label{frames:char:trace-subsection}

We use a unitary realization of each irreducible type of a symmetric
group. For type $\alpha$, write $D_\alpha$ for its degree and
$\chi_\alpha$ for its character. We use the ordinary,
unnormalized Hilbert--Schmidt norm. For $m\ge2$ and $u\in S_m$, let
$c_i(u)$ be the number of $i$-cycles and define $e_1,\ldots,e_m$ by
\[
 \sum_{i=1}^k e_i=
 \frac{\log\max\{1,\sum_{i=1}^k i c_i(u)\}}{\log m}
 \quad(1\le k\le m),\qquad
 E(u)=\sum_{i=1}^m\frac{e_i}{i}.
\]
The $e_i$ are nonnegative and sum to one. Larsen and Shalev's theorem
states that, for every $\epsilon>0$, for all sufficiently large $m$
depending only on $\epsilon$, and uniformly over $u$ and $\alpha$,
\begin{equation}\label{frames:char:LS}
 |\chi_\alpha(u)|\le D_\alpha^{E(u)+\epsilon}.
\end{equation}
See \cite[Theorem~1.1]{larsen-shalev2008} and the formulation with
fixed slack in
\cite[Definition~2.1 and Theorem~2.2]{keller-lifshitz-sheinfeld2024}.
If $u$ has at most $m^{9/10}$ fixed points, then
$E(u)\le e_1+(1-e_1)/2\le19/20$. Taking $\epsilon=3/100$ gives
\begin{equation}\label{frames:char:few-fixed}
 \frac{|\chi_\alpha(u)|}{D_\alpha}\le D_\alpha^{-1/50}.
\end{equation}
This is a pointwise bound and requires no centrality of a shuffle law.

For every fixed $s>0$, the Witten-zeta estimate is
\begin{equation}\label{frames:char:zeta}
 \sum_{\alpha\in\widehat{S_m}}D_\alpha^{-s}=2+o(1)
 \quad(m\to\infty).
\end{equation}
See \cite[Theorems~1.1 and~2.6]{liebeck-shalev2004} and
\cite[Theorem~2.14]{keller-lifshitz-sheinfeld2024}.
In particular the complete sum at $s=1$ is bounded by an absolute
$C_\zeta$ for every $m\ge1$, and the sum with the trivial and sign types
removed tends to zero. Lemma~\Lref{l:degree-inverse-first} of
\cite{partial-fourier} supplies an elementary proof of these last facts.

\begin{lemma}\label{frames:char:one-pass}
Let $b=2^h$ and let $Y$ apply each of the $h$ coordinate layers once,
in any order, using independent fair switches. For every irreducible
representation $\rho_\alpha$ of $S_b$, put $T_\alpha=\mathbb E\rho_\alpha(Y)$.
Uniformly for all sufficiently large $b$,
\[
 \frac{\operatorname{tr}(T_\alpha^*T_\alpha)}{D_\alpha}
 =\frac{\operatorname{tr}(T_\alpha T_\alpha^*)}{D_\alpha}
 \le2D_\alpha^{-1/100}.
\]
\end{lemma}

\begin{proof}
Let $Y'$ be an independent copy and put $u=(Y')^{-1}Y$. Independence
and unitarity give
\[
 \frac{\operatorname{tr}(T_\alpha^*T_\alpha)}{D_\alpha}
 =\mathbb E\frac{\chi_\alpha(u)}{D_\alpha}
 \le\mathbb E\frac{|\chi_\alpha(u)|}{D_\alpha}.
\]
If $f$ counts the fixed points of $u$, then, for $1\le k\le b$,
\begin{equation}\label{frames:char:fixed-tail}
 \Pr(f\ge k)\le\binom bk b^{-k/2}.
\end{equation}
Indeed, condition on $Y'$ and prescribe $Y(x)=Y'(x)$ at a chosen set
of $k$ inputs. Each input-output pair determines one route because each
coordinate is updated once. A consistent family of routes visits at
least $hk/2$ distinct switch bits: a switch serves at most two routes.
It therefore has probability at most $2^{-hk/2}=b^{-k/2}$. A union
bound over the input set proves \eqref{frames:char:fixed-tail}.

Write $D=D_\alpha$. The assertion is immediate for $D=1$. For $D>1$ set
\[
 \kappa=\frac{\log D}{2\log b},\qquad
 k=\left\lfloor\max\{b^{9/10},\kappa\}\right\rfloor+1.
\]
Since $D^2\le b!$, one has $\kappa\le b/4$, so $k\le b$ for large $b$.
For $k\ge b^{9/10}$ and sufficiently large $b$,
\[
 \binom bk b^{-k/2}
 \le\left(\frac{e b^{1/2}}k\right)^k
 \le b^{-k/4}\le D^{-1/8}.
\]
Outside an event of probability at most $D^{-1/8}$, therefore,
$f\le b^{9/10}$ or $f\le\kappa$. The first case is covered by
\eqref{frames:char:few-fixed}. In the second, restrict $\rho_\alpha$
to the symmetric group on the $b-f$ moving points. Every irreducible
constituent of degree $E$ satisfies
\[
 E\ge\frac{D}{[S_b:S_{b-f}]}\ge b^{-f}D\ge D^{1/2}.
\]
For the first inequality, translates of one constituent subspace by
coset representatives span a nonzero $S_b$-invariant subspace and hence
the whole irreducible representation. Since $f\le\kappa\le b/4$,
the moving-point group has size tending uniformly to infinity.
The permutation has no fixed points there. Applying
\eqref{frames:char:few-fixed} to each constituent and averaging with
its dimension weight, including multiplicity, bounds the normalized
character by $D^{-1/100}$. On the exceptional event use the bound one.
Thus the normalized trace is at most
$D^{-1/100}+D^{-1/8}\le2D^{-1/100}$.
\end{proof}

\paragraph{An alternative restriction estimate.}
The preceding proof lower-bounds every constituent degree after fixed
points are removed. A distinct Young-branching argument controls the
number of constituents and gives another estimate for the same trace.

\begin{lemma}\label{char:sweep-trace}
Let $s=2^\ell$ and let $G_s$ be a sweep using each of the $\ell$ coordinate
directions once, in any order, with independent fair switches. There are
absolute $a>0,C_1<\infty$ such that, for every sufficiently large $s$ and
every complex irreducible representation $\rho_\alpha$ of $S_s$,
\[
 q_\alpha:=D_\alpha^{-1}
  \left\|\mathbb E\rho_\alpha(G_s)\right\|_{\mathrm{HS}}^2
 \le C_1D_\alpha^{-a}.
\]
The Hilbert--Schmidt norm uses the ordinary, unnormalized trace.
\end{lemma}

\begin{proof}
For a permutation of $m$ points with at most $m^{4/5}$ fixed points,
one has $e_1\le4/5$ and hence $E(g)\le9/10$.
Applying \eqref{frames:char:LS} with slack $1/20$ yields
\begin{equation}\label{char:few-fixed}
 |\chi_\alpha(g)|\le D_\alpha^{19/20}
\end{equation}
for all sufficiently large $m$, uniformly in $g$ and $\alpha$.

Let $G_s'$ be an independent copy, put $u=(G_s')^{-1}G_s$, let $f$
count its fixed points, and write $D=D_\alpha$. As in the preceding
trace calculation, $q_\alpha=\mathbb E\chi_\alpha(u)/D$.
The case $D=1$ is immediate. When $f\le s^{4/5}$,
\eqref{char:few-fixed} bounds the absolute normalized character by
$D^{-1/20}$. Suppose instead that
\[
 s^{4/5}<f\le\kappa:=\frac{\log D}{2\log s}.
\]
Since $D\le s!\le s^s$, one has $\kappa\le s/2$ and hence $f\le s/2$
in this case. The moving-point group
therefore has size tending uniformly to infinity. Iterated Young
branching \cite[Theorem~9.2]{James1978} expresses the restriction to
that group as at most $s^f$ irreducible summands counted with
multiplicity. If their degrees are $E_1,\ldots,E_r$, then
$\sum_jE_j=D$. The permutation is fixed-point-free on this group.
Applying \eqref{char:few-fixed} there and using concavity gives
\[
 |\chi_\alpha(u)|
 \le\sum_{j=1}^r E_j^{19/20}
 \le r^{1/20}D^{19/20}
 \le s^{f/20}D^{19/20}
 \le D^{39/40}.
\]
For the remaining event take
$k=\lfloor\max\{s^{4/5},\kappa\}\rfloor+1$, which lies in $[1,s]$
for large $s$. The unique-route tail
\eqref{frames:char:fixed-tail}, with $s$ in place of $b$, gives for
large $s$
\[
 \Pr(f\ge k)\le\binom sk s^{-k/2}
 \le(e s^{1/2}/k)^k\le s^{-k/5}\le D^{-1/10}.
\]
The normalized trace is consequently at most
$D^{-1/40}+D^{-1/10}\le2D^{-1/40}$. This proves the stated existence
with the distinct branching conclusion $a=1/40$ and $C_1=2$.
\end{proof}

The following two conversions use the shared bound with exponent
$1/100$. The alternative above gives the stronger trace exponent
$1/40$ through the different restriction argument.

\subsection{Smoothing after the mixing segment: capped preimages}
\label{frames:char:smoothing}

For the first block order, take $L=100d$ layers with product $\pi$,
followed by $d-1$ layers with product $\sigma$. The last layers omit
one coordinate and preserve its two halves $J_0,J_1$, each of size
$b=n/2$. Their fair coins give independent one-pass permutations
$\sigma_0,\sigma_1$ on these halves, independent also of $\pi$.
A block has length $101d-1$. Its phase and its missing axis may change
from one block to the next.

Fix an ordering of each $J_j$, and let $Q_j(\pi)$ be its ordered
preimage list. Lemma~\ref{frames:char:half-list} gives a bound
$\epsilon_n=o(1)$ for the distance of each list law $\nu_j$ from the
uniform injection law $U_b$, uniformly in phase. Delete every $\pi$
for which either list has density greater than two, and condition on
what remains. If $B_j$ is the excluded set of lists, then
\[
 \nu_j(B_j)\le2\bigl(\nu_j(B_j)-U_b(B_j)\bigr)\le2\epsilon_n.
\]
The retained event has probability at least $1-4\epsilon_n$. For large
$n$ it has probability at least one half, so each separate preimage
marginal in the modified law is bounded by $4U_b$. The modification
costs at most $4\epsilon_n$ in total variation. The independent passes
are unchanged.

\begin{lemma}\label{frames:char:minimum-degree}
For even $n\ge6$, every irreducible representation of $S_n$ other than
trivial and sign has degree at least $n/2$.
\end{lemma}
\begin{proof}
Restrict to the elementary abelian subgroup generated by a fixed matching
of $n/2$ transpositions. If every constituent assigns the same sign to
all these generators, their representing matrices are equal. Conjugation
then makes any two disjoint transpositions have equal matrices. A third
transposition disjoint from two overlapping ones exists when $n\ge6$,
so all transpositions have the same matrix. The image is generated by
one involution, and irreducibility gives the trivial or sign representation.
Otherwise a constituent assigns negative sign to exactly $r$ generators,
with $1\le r\le n/2-1$. The normalizer permutes the matching pairs,
so all $\binom{n/2}{r}\ge n/2$ reordered characters occur. Their
eigenspaces are linearly independent.
\end{proof}

\begin{proposition}\label{frames:char:block-norm}
There is an absolute $C_0$ such that, for every sufficiently large
$n=2^d$, every modified block above, and every irreducible $\rho$ of
$S_n$ of degree $D$ other than trivial or sign,
\[
 \|\mathbb E\rho(\sigma\pi)\|_{\mathrm{op}}
 \le C_0D^{-1/20000}.
\]
The modified block has zero sign coefficient. These statements are
uniform in its starting phase.
\end{proposition}

\begin{proof}
Let $G_j$ be the symmetric group supported on $J_j$ and put
$H=G_0\times G_1$. If a positive operator $Q$ commutes with
$\rho(G_{1-j})$, then $\rho(\pi)^*Q\rho(\pi)$ is determined by
$Q_j(\pi)$. Indeed, equal ordered preimages of $J_j$ mean
$\pi'\pi^{-1}\in G_{1-j}$. The separate marginal cap and positivity
therefore give, for every unit vector $v$,
\begin{equation}\label{frames:char:schur-cap}
 \mathbb E\langle\rho(\pi)v,Q\rho(\pi)v\rangle
 \le4\,\mathbb E_{g\sim U_{S_n}}
       \langle\rho(g)v,Q\rho(g)v\rangle
 =4\,\frac{\operatorname{tr}Q}{D}.
\end{equation}
The final equality is Schur averaging. This comparison uses one capped
preimage marginal at a time; it imposes no joint density bound.

Set $A_j=\mathbb E\rho(\sigma_j)$ and $A=A_0A_1$. These factors and
their adjoints commute, and they are contractions. Independence and
convexity give
\begin{equation}\label{frames:char:convexity}
 \|A\,\mathbb E\rho(\pi)v\|^2
 \le\mathbb E\langle\rho(\pi)v,A^*A\rho(\pi)v\rangle.
\end{equation}
For each $j$, let $P_{j,\le}$ select the full $G_j$-isotypic spaces
of degree at most $D^{1/100}$, and put $P_{j,>}=I-P_{j,\le}$.
These projections commute with $H$ and with each other. On a joint
half-type, write $X_j=A_j^*A_j$. The $X_j$ are positive contractions
on separate tensor factors. If type $j$ is large, $X_0X_1\preceq X_j$;
if both types are small, $X_0X_1\preceq I$. Hence
\begin{equation}\label{frames:char:positive-comparison}
 A^*A\preceq
 A_0^*A_0P_{0,>}+A_1^*A_1P_{1,>}+P_{0,\le}P_{1,\le}.
\end{equation}
Every term is positive. The first two commute with the complementary
half group, and the last commutes with both groups, so
\eqref{frames:char:schur-cap} applies to each.

On each $G_j$-type selected by $P_{j,>}$,
Lemma~\ref{frames:char:one-pass} bounds
the normalized trace by
$2D_\alpha^{-1/100}\le2D^{-1/10000}$. Averaging with the restriction
dimension weights, including all multiplicities, gives
\begin{equation}\label{frames:char:large-trace}
 \frac1D\operatorname{tr}(A_j^*A_jP_{j,>})
 \le2D^{-1/10000}.
\end{equation}

It remains to bound the fraction of the space on which both half-types
are small. If $m_{\alpha\beta}$ is the multiplicity of
$\alpha\otimes\beta$ in $\rho|_H$, then
\begin{equation}\label{frames:char:multiplicity}
 m_{\alpha\beta}\le\sqrt{b+1}\,D_\alpha D_\beta.
\end{equation}
By Frobenius reciprocity \cite[Theorem~4.33]{etingof},
$m_{\alpha\beta}$ is the multiplicity of $\rho$ in
$W=\operatorname{Ind}_H^{S_n}(\alpha\otimes\beta)$, so
$m_{\alpha\beta}^2\le\dim\operatorname{End}_{S_n}(W)$.
Realize $W$ in fibers of dimension $D_\alpha D_\beta$ over $S_n/H$.
An equivariant block map at one representative ordered pair of cosets
determines its transported maps on that pair orbit. The stabilizer can
impose linear constraints on the representative map and can only reduce
the number of choices. Cosets correspond to $b$-subsets, and intersection
size classifies their ordered-pair orbits, of which there are $b+1$.
Each representative map has at most $(D_\alpha D_\beta)^2$ parameters.
This proves \eqref{frames:char:multiplicity}.

The complete inverse-first sum gives
\[
 \sum_{D_\alpha\le D^{1/100}}D_\alpha^2
 \le D^{3/100}\sum_{\alpha\in\widehat{S_b}}D_\alpha^{-1}
 \le C_\zeta D^{3/100}.
\]
Lemma~\ref{frames:char:minimum-degree} gives $b\le D$. Therefore
\begin{equation}\label{frames:char:small-trace}
 \frac{\operatorname{tr}(P_{0,\le}P_{1,\le})}{D}
 \le\frac{\sqrt{b+1}}D
       \left(\sum_{D_\alpha\le D^{1/100}}D_\alpha^2\right)^2
 =O(D^{-1/2+6/100}).
\end{equation}
Equations~\eqref{frames:char:schur-cap}--\eqref{frames:char:small-trace}
bound the squared operator norm by
$16D^{-1/10000}+O(D^{-1/2+6/100})$. Taking square roots proves the
stated exponent. Finally, the pass $\sigma$ contains an independent
fair switch for large $d$. Toggling it changes sign, and its independence
from the modified $\pi$ gives zero sign coefficient for the block.
\end{proof}

\subsection{Completion of the capped-preimage route}
\label{frames:char:amplification}

\begin{proof}[Proof of Theorem~\ref{frames:char:main}]
Take $r=60000$ successive blocks, modifying their long segments
independently. The product law $\mu$ differs from the unmodified product
by at most $r\,o(1)=o(1)$. The block phases may differ, but the estimates
are uniform. Their Fourier matrices multiply in chronological order,
and Proposition~\ref{frames:char:block-norm} gives
\[
 \|\mathbb E_\mu\rho\|_{\mathrm{op}}
 \le C_1D_\rho^{-3},
 \qquad C_1=C_0^{60000},
 \qquad \rho\notin\{\mathbf1,\operatorname{sign}\}.
\]
The sign coefficient is zero. Plancherel, the Hilbert--Schmidt inequality,
and Lemma~\ref{frames:char:minimum-degree} yield
\[
 4\|\mu-U_{S_n}\|_{\mathrm{TV}}^2
 \le C_1^2\sum_{\rho\notin\{\mathbf1,\operatorname{sign}\}}D_\rho^{-4}
 \le C_1^2(2/n)^3\sum_{\rho\in\widehat{S_n}}D_\rho^{-1}=o(1).
\]
Restoring the modified segments still costs $o(1)$.

The unmodified coordinate product has
$t=60000(101d-1)$ layers and equals $R^{-t}Y_t$ in the notation of
Section~\ref{sec:frames}. Restoring the one terminal rotation preserves
uniform measure and total variation even when $t$ is not divisible by
$d$. A deterministic initial deck is a right translate, so it has the
same distance. Since $t\le6060000d$, the upper bound follows. The lower
bound is the earlier one-card or support lower bound.
\end{proof}
